
\section{Evaluation}
\subsection{Experimental Setup}
\label{subsec:setup}

\subsubsection{Hardware and Software Environment}

All experiments run on a single-node Spark deployment, with the data residing
permanently in object storage and no local replica produced during query
execution. Both evaluation environments co-locate compute and storage in the
same region. This configuration minimizes network round-trip latency and
therefore constitutes the most conservative setting for evaluating read-path
optimizations: under same-region deployment, the fraction of wall-clock time
attributable to remote reads is at its smallest, so the gains reported in this
paper constitute a lower bound on the gains these mechanisms deliver in
higher-latency environments.

The two environments are deployed on AWS and Tencent Cloud, respectively. The
AWS environment uses an EC2 \texttt{m5dn.4xlarge} instance (16 vCPUs, 62\,GiB
of memory) with S3 \texttt{us-east-2} as object storage; the compute engine is
Spark~4.1 running \texttt{local[16]} with a driver heap (\texttt{Xmx}) of
48\,GiB; the Parquet stack is parquet-mr~1.15.x with vectored I/O enabled;
Spark scratch space is placed on a dedicated instance-store NVMe device,
physically separated from the cache medium; garbage collection uses G1 with a
periodic GC interval of 30\,s; and the cache tier is an on-heap 4\,GiB cache.
The Tencent Cloud environment uses a Tencent Cloud CVM instance with 32 vCPUs
and COS \texttt{ap-guangzhou} as object storage, and is otherwise identical to
the AWS environment. The concurrency and memory conditions on the execution
side are thus comparable across the two environments: although their physical
core counts are 16 and 32, \texttt{local[16]} fixes the effective degree of
parallelism at 16 in both cases. The cache capacity is fixed at 4\,GiB; the
rationale for this value follows from the capacity curves in
\S\ref{subsec:capacity}.

\subsubsection{Datasets}

We select five datasets covering two classes of read workloads, listed in
Table~\ref{tab:datasets}. The OLAP workloads are taken from two standard
benchmarks, which provide comparability and reproducibility, whereas the
feature-scan workloads are taken from three real machine-learning datasets,
which introduce the access pattern in which \emph{the same set of columns is
repeatedly scanned in full}---a pattern characteristic of training pipelines.
Among the two OLAP workloads, ClickBench is a single wide table exercised by
highly selective predicates and aggregations with strongly skewed column
access, while TPC-H SF350 involves eight-table joins and exhibits both the
highest degree of object fragmentation and the largest volume of remote reads.
Among the three feature-scan workloads, NYC Yellow Taxi repeats four epochs and
exhibits the strongest temporal locality, Criteo CTR comprises the largest
number of objects, and NovoMCP Compound comprises the fewest objects with the
largest individual object size.

This collection spans a wide range along three axes that bear directly on the
read path: the number of objects in the canonical layout ranges from 21 for
NovoMCP to 2{,}130 for TPC-H; the ratio of the per-query remote footprint to
the cache capacity ranges from comparable to the capacity to roughly six times
larger; and the reuse structure falls into two classes, namely explicit
repetition and column overlap. Both the effective and the ineffective regime of
each mechanism therefore fall within the sample.

Beyond the SF8 configuration used throughout this section, ClickBench is
additionally measured at SF1 (99{,}997{,}497 rows, 14.11\,GiB) exclusively for
the scale-sensitivity analysis of \S\ref{subsec:scale}; that configuration is
excluded from the grids reported below.

% requires: booktabs, makecell
% Preamble:
% \usepackage{booktabs}
% \usepackage{makecell}
% \usepackage{tabularx}
% \usepackage{array}

\begin{table}[t]
\centering
\scriptsize
\setlength{\tabcolsep}{3pt}
\caption{Datasets used in the evaluation.}
\label{tab:datasets}

\begin{tabularx}{\columnwidth}{@{}
  >{\raggedright\arraybackslash}p{0.17\columnwidth}
  >{\raggedright\arraybackslash}p{0.12\columnwidth}
  >{\raggedleft\arraybackslash}p{0.17\columnwidth}
  >{\raggedright\arraybackslash}X
  >{\raggedright\arraybackslash}X
  >{\centering\arraybackslash}p{0.08\columnwidth}
@{}}
\toprule
\textbf{Dataset} &
\makecell[l]{\textbf{Work-}\\\textbf{load}} &
\textbf{Rows} &
\makecell[l]{\textbf{Canonical}\\\textbf{layout}} &
\makecell[l]{\textbf{Reconfig}\\\textbf{layout}} &
\makecell[c]{\textbf{Que-}\\\textbf{ries}} \\
\midrule

\makecell[l]{ClickBench\\SF8} &
OLAP &
799{,}979{,}976 &
\makecell[l]{881 obj. /\\111.99\,GiB} &
\makecell[l]{176 obj. /\\81.92\,GiB} &
43 \\

\addlinespace[1pt]

\makecell[l]{TPC-H\\SF350} &
OLAP &
3{,}031{,}425{,}280$^{\dagger}$ &
\makecell[l]{2{,}130 obj. /\\117.48\,GiB} &
\makecell[l]{514 obj. /\\85.53\,GiB} &
22 \\

\addlinespace[1pt]

\makecell[l]{NYC Yellow\\Taxi} &
\makecell[l]{Feature\\scan} &
1{,}826{,}960{,}642 &
\makecell[l]{204 obj. /\\39.28\,GiB} &
\makecell[l]{109 obj. /\\38.39\,GiB} &
4$^{\ddagger}$ \\

\addlinespace[1pt]

\makecell[l]{Criteo\\CTR} &
\makecell[l]{Feature\\scan} &
747{,}844{,}676 &
\makecell[l]{973 obj. /\\43.61\,GiB} &
\makecell[l]{32 obj. /\\42.02\,GiB} &
24 \\

\addlinespace[1pt]

\makecell[l]{NovoMCP\\Compound} &
\makecell[l]{Feature\\scan} &
122{,}454{,}458 &
\makecell[l]{21 obj. /\\34.27\,GiB} &
\makecell[l]{20 obj. /\\29.07\,GiB} &
20 \\

\bottomrule
\end{tabularx}

\vspace{2pt}
\begin{minipage}{\columnwidth}\scriptsize
$^{\dagger}$Summed over all eight tables and derived from the SF350
cardinalities in the TPC-H specification, which defines this cardinality only
approximately. All other row counts are measured.
\quad
$^{\ddagger}$The four queries are four epochs of the same scan, differing only
in an epoch literal that keeps the statement texts distinct.
\end{minipage}
\end{table}


\subsubsection{Protocol and Latency Characterization}

\textbf{$2\times2$ factorial design.}
We treat each of the two mechanisms as a binary factor: layout $\in$
\{canonical, reconfigured\} and cache $\in$ \{off, on\}, yielding four cells
denoted 00, 10, 01, and 11. The purpose of the factorial design is to obtain
both main effects together with their interaction term, the latter
characterizing the influence of one mechanism on the effectiveness of the other
(see \S\ref{subsec:interaction}). The interaction term is defined as
$$
  \Delta = (01 - 11) - (00 - 10),
$$
that is, the difference between the time the cache saves on the reconfigured
layout and the time it saves on the canonical layout, which is identical to the
combined gain minus the sum of the two individual gains. $\Delta > 0$ indicates
that the two mechanisms reinforce each other, whereas $\Delta < 0$ indicates
that their gains overlap.

\textbf{Control of run-to-run variation.}
Each cell is run for $n=5$ rounds, and the median of the per-round sum of
query wall-clock times is taken as the representative value for that cell. So
that drift in machine state remains separable from mechanism effects, the four
cells are executed in an interleaved fashion within each round, which
distributes slow drift evenly across cells; each cell runs in an independently
cold-started JVM with a cold cache.

\textbf{Cross-cloud comparison.}
The same data snapshot, the same build artifact, and the same query set are
used to run one complete grid in each of the two cloud environments. This
design introduces \emph{the time value of a single eliminated request} as a
controlled variable: the physical reduction (in remote request count and bytes)
is determined jointly by the layout and the cache and is preserved across the
two environments, whereas the same reduction translates into different
wall-clock gains under different per-request costs. Criteo provides the
clearest instance of this effect---the I/O metrics of the two environments
agree almost digit for digit (the reconfigured layout issues 44{,}242 GETs and
reads 187.87\,GiB in both; with caching enabled, 15{,}357 and 15{,}510 GETs,
respectively)---so that the difference in wall-clock time is attributable to
the difference in per-request cost.

The fixed per-request costs of the two environments differ by a factor of
approximately $3.3$, as shown in Table~\ref{tab:latency}. This observation
constitutes the direct rationale for both mechanisms in this paper taking the
reduction of request counts as their primary means.

\begin{table}[t]
\centering
\footnotesize
\caption{Per-request latency of the two object stores (same-region).}
\label{tab:latency}
\begin{tabular}{@{}lrr@{}}
\toprule
Request size & AWS S3 (p50) & Tencent COS (p50) \\
\midrule
Small GET ($<$ 64\,KiB) & 23.6\,ms & 81\,ms \\
Large GET ($\geq$ 8\,MiB) & 48.3\,ms & 112\,ms \\
Fitted per-request cost & 26.1\,ms & 86.2\,ms \\
\bottomrule
\end{tabular}
\end{table}

\subsection{Main Results}
\label{subsec:main}

This section reports the end-to-end gains obtained when both mechanisms are
enabled jointly. Table~\ref{tab:main} presents the combined effect across all
ten grids: within each grid, the configuration with both mechanisms enabled
(cell~11) is compared against that grid's own canonical baseline (cell~00).

\textbf{The combined gain holds across all grids.}
The median wall-clock reduction over the ten grids is 31.3\%, with values
ranging from 12.0\% (NovoMCP / AWS) to 69.7\% (NYC Yellow Taxi / COS). Within
every grid the per-round paired differences are consistent in sign, so the gain
holds under the significance criterion adopted in this paper. Separated by
cloud, the median reduction is 24.6\% on AWS and 49.9\% on COS. The gain holds
at both extremes of workload scale: the largest workload, TPC-H SF350 / COS
(canonical wall clock 10{,}709\,s), attains a reduction of 52.2\%, while the
smallest, NovoMCP / AWS (canonical wall clock 197.0\,s), attains 12.0\%. The
mechanisms therefore neither lose effectiveness at large scale nor have their
gains absorbed by fixed overheads at small scale. Figure~\ref{fig:combined}
shows the per-grid reductions with the two clouds placed side by side.

\begin{table}[t]
\centering
\caption{End-to-end results per cloud. Row order is identical in both
subtables. The second column gives the canonical baseline (cell~00) and the
combined configuration (cell~11); wall-clock reductions are relative to
cell~00 of the same grid.}
\label{tab:main}
\scriptsize
\setlength{\tabcolsep}{3pt}
\renewcommand{\arraystretch}{1.15}

\textbf{(a) AWS S3 (\texttt{us-east-2})}\\[2pt]

\resizebox{\columnwidth}{!}{%
\begin{tabular}{@{}lrrr@{}}
\toprule
Dataset &
00 (s) / 11 (s) &
Wall clock reduction &
Peak heap 00 / 11 (GiB) \\
\midrule
ClickBench SF8 &
2{,}395.5 / 1{,}638.9 &
$-31.6\%$ &
3.56 / 23.60 \\

\addlinespace[1pt]
TPC-H SF350 &
6{,}775.9 / 4{,}676.9 &
$-31.0\%$ &
14.45 / 43.07 \\

\addlinespace[1pt]
NYC Yellow Taxi &
345.3 / 275.0 &
$-20.4\%$ &
1.99 / 17.19 \\

\addlinespace[1pt]
Criteo CTR &
1{,}774.4 / 1{,}338.2 &
$-24.6\%$ &
11.97 / 15.12 \\

\addlinespace[1pt]
NovoMCP Compound &
197.0 / 173.4 &
$-12.0\%$ &
13.25 / 21.59 \\

\midrule
Median &
&
$-24.6\%$ &
\\
\bottomrule
\end{tabular}%
}

\vspace{6pt}

\textbf{(b) Tencent Cloud COS (\texttt{ap-guangzhou})}\\[2pt]

\resizebox{\columnwidth}{!}{%
\begin{tabular}{@{}lrrr@{}}
\toprule
Dataset &
00 (s) / 11 (s) &
Wall clock reduction &
Peak heap 00 / 11 (GiB) \\
\midrule
ClickBench SF8 &
2{,}890.3 / 1{,}446.6 &
$-49.9\%$ &
9.10 / 22.00 \\

\addlinespace[1pt]
TPC-H SF350 &
10{,}709.0 / 5{,}115.2 &
$-52.2\%$ &
22.00 / 42.90 \\

\addlinespace[1pt]
NYC Yellow Taxi &
632.0 / 191.7 &
$-69.7\%$ &
2.10 / 16.78 \\

\addlinespace[1pt]
Criteo CTR &
1{,}357.5 / 947.8 &
$-30.2\%$ &
7.64 / 20.64 \\

\addlinespace[1pt]
NovoMCP Compound &
257.4 / 158.1 &
$-38.6\%$ &
9.82 / 20.54 \\

\midrule
Median &
&
$-49.9\%$ &
\\
\bottomrule
\end{tabular}%
}
\end{table}

\begin{figure}[htbp]
    \centering
    \includegraphics[width=0.5\textwidth]{figs/recas_combined.png}
    \caption{Combined wall-clock reduction across the ten grids, with AWS and
  COS placed side by side.}
    \label{fig:combined}
\end{figure}

\textbf{The gain originates in request count rather than in bytes.}
The reduction in remote request count is highly consistent across the grids,
falling between 81.0\% and 95.6\% in all ten of them (median 83.5\%, a spread
of 15 percentage points), whereas the reduction in remote bytes spans 3.5\% to
43.4\% (a spread of 40 percentage points) and bears no correspondence to the
wall-clock reduction. Two extreme samples illustrate this most clearly: NYC
Yellow Taxi / COS attains the largest wall-clock reduction in the entire sample,
69.7\%, from a byte reduction of only 3.6\%; and TPC-H SF350 / AWS likewise
reduces bytes by only 6.0\% while realizing a reduction of 31.0\%.

This result corroborates the latency characterization of
\S\ref{subsec:protocol}: in the same-region AWS environment the fitted cost of
a single GET is 26.1\,ms, of which approximately 24\,ms is a server-side
time-to-first-byte floor that is independent of request size and does not
improve with bandwidth. Wall-clock time on a lakehouse read path is therefore
governed by the \emph{number} of remote requests rather than by their
\emph{volume}, and the two mechanisms act on precisely this common surface:
layout reconfiguration reduces the requests that are issued in the first place
by coalescing physical reads, while caching eliminates the repeated issuing of
requests by reusing bytes already fetched.

\subsubsection{Resource Cost}
\label{subsec:cost}

The gains above are obtained at the cost of memory occupancy; the last column
of Table~\ref{tab:main} gives the peak heap of cells~00 and~11 for each grid.

\textbf{The heap is the dominant cost.}
With both mechanisms enabled, the increase in peak heap ranges from 3.15\,GiB
(Criteo / AWS) to 28.62\,GiB (TPC-H SF350 / AWS); the largest absolute values
occur on TPC-H SF350, reaching 43.07 and 42.90\,GiB on the two clouds,
approximately 90\% of the 48\,GiB heap ceiling. This increase substantially
exceeds the 4\,GiB cache budget, because it is not composed solely of cached
bytes: SDK and Spark objects kept alive along the hit path, the misalignment of
peak instants across query phases, and garbage-collection timing all contribute
to the observed value. The peak-heap difference is thus the total memory cost of
enabling the cache rather than the capacity of the cache itself. This fact
determines that the scalability ceiling of an in-memory cache tier is set by
the heap rather than by the hit rate, and \S\ref{subsec:capacity} fixes the
4\,GiB capacity on that basis.

\textbf{Garbage-collection overhead does not rise with the heap.}
While heap occupancy increases substantially, garbage-collection time remains
unchanged or decreases across the grids: it falls from 30.9\,s to 17.2\,s on
ClickBench SF8 / AWS and from 77.9\,s to 13.1\,s on Criteo / AWS, and stays
uniformly between 4.0 and 5.4\,s across the four cells of NYC Yellow Taxi. The
mechanism is that reducing the request count simultaneously reduces short-lived
allocations: metadata objects and read buffers are produced in proportion to the
number of requests, so once the request count falls by an order of magnitude,
this allocation pressure disappears and offsets the collection cost incurred by
long-lived cached blocks. The memory cost therefore manifests as occupancy of
resident capacity rather than as additional collection time.

\textbf{Layout reconfiguration affects the heap in both directions.}
When layout reconfiguration is enabled alone, the peak heap decreases on some
workloads and increases on others, depending on the bottleneck of the workload
in question. It falls from 11.97 to 4.78\,GiB on Criteo / AWS and from 13.25 to
6.50\,GiB on NovoMCP / AWS, because contracting the file count from 973 to 32
and from 21 to 20, respectively, sharply lowers the short-lived allocation of
metadata and read objects. It rises from 1.99 to 12.89\,GiB on NYC Yellow Taxi
/ AWS, because reconfiguration markedly enlarges individual reads (109.33\,GiB
carried by only 9{,}266 GETs), so that more decode and aggregation batches are
alive concurrently. Both effects follow from the change in request shape, and
their direction is determined by the degree of fragmentation originally present
in the workload.

\subsection{Sensitivity Analysis: The Conditions Under Which Each Mechanism Is Effective}
\label{subsec:sensitivity}

\S\ref{subsec:main} reported the combined gain. This section decomposes it into
the individual effects of the two mechanisms and establishes that both
optimizations in our design are necessary and cannot substitute for each other:
their effectiveness is governed by two mutually independent properties of the
workload, so that omitting either one leaves part of the workload space
uncovered. We then examine the moderating influence of two conditioning
variables---environment latency and data scale---and finally characterize the
interaction between the two mechanisms. Figure~\ref{fig:decomposition} gives the
decomposition into the Caching, Layout, and Combined effects for each grid.

\subsubsection{Effectiveness of Caching: Per-Query Remote Read Volume Relative to Capacity}
\label{subsec:pressure}

The effectiveness of caching is governed by the ratio between the volume of
remote data read by a single query and the cache capacity. Starting from the
request-level I/O trace of the canonical baseline (cell~00), we attribute every
remote request to its originating query and thereby obtain the \emph{per-query
remote read volume}; the ratio of this quantity to the 4\,GiB cache capacity is
termed the \emph{capacity pressure ratio}. Because queries differ greatly in
weight---the per-query read volume of ClickBench SF8 spans 0.04 to
111.66\,GiB, three orders of magnitude---we take the median as the
representative value of the distribution.

\begin{table}[t]
\centering
\caption{Capacity pressure ratio and the request reduction achieved by caching
alone. Pressure ratio is the median per-query read volume divided by the
4\,GiB cache capacity.}
\label{tab:pressure}
\footnotesize
\setlength{\tabcolsep}{4pt}
\begin{tabular}{@{}lrrr@{}}
\toprule
Dataset & \makecell[r]{Median per-query\\read volume} &
\makecell[r]{Pressure\\ratio} & \makecell[r]{GET reduction\\by caching} \\
\midrule
ClickBench SF8   & 2.71\,GiB  & 0.7$\times$ & 73.0\% \\
NovoMCP Compound & 2.74\,GiB  & 0.7$\times$ & 77.2\% \\
Criteo CTR       & 4.92\,GiB  & 1.2$\times$ & 74.8\% \\
NYC Yellow Taxi  & 26.00\,GiB & 6.5$\times$ & \textbf{10.6\%} \\
TPC-H SF350      & 26.81\,GiB & 6.7$\times$ & \textbf{11.3\%} \\
\bottomrule
\end{tabular}
\end{table}

The relationship exhibits a threshold form: the three datasets whose pressure
ratio does not exceed 1.2 see 73.0\% to 77.2\% of their requests eliminated,
whereas the two datasets with a ratio of 6.5 or above see the reduction collapse
to roughly 11\%. ClickBench SF8 and NovoMCP have almost identical pressure
ratios (0.68 and 0.69) and correspondingly similar request reductions (73.0\%
and 77.2\%), which constitutes an independent corroboration of the metric.

A more direct mechanistic characterization is the \emph{fraction of queries
whose read volume does not exceed the cache capacity}. Of the 43 queries in
ClickBench SF8, 26 (60\%) read less than 4\,GiB; of the 22 queries in TPC-H
SF350, only one does; and of the four in NYC Yellow Taxi, none does. The
request reduction achieved by caching varies in the same direction as this
fraction, indicating that the gain does not accrue uniformly from all queries
but is concentrated in those whose read volume falls within the capacity.
ClickBench is the canonical illustration: its large scans---Q21 through Q24,
Q28, Q29, Q34, and Q35---each read between 17 and 112\,GiB and are refused
admission by the frequency filter, while the 73.0\% request reduction arises
from the repeated access of the remaining small queries to the same set of hot
columns.

NYC Yellow Taxi provides a counter-intuitive instance. Its temporal locality is
the strongest among the five workloads---its query set consists entirely of four
epochs of one identical scan, so that reuse is deterministic rather than
incidental. The four epochs nevertheless read 25.96, 25.98, 26.01, and
26.14\,GiB respectively, none of which fits within the 4\,GiB capacity, so the
cache is flushed in its entirety before the next epoch begins and only 10.6\% of
requests are eliminated. \textbf{The effectiveness of caching therefore does not
follow from temporal locality alone; it is additionally governed by the capacity
pressure ratio.}

On workloads whose pressure ratio far exceeds unity, the cache degenerates into
a small-object cache: the 11.3\% reduction on TPC-H comes from the
\texttt{nation}, \texttt{region}, and \texttt{supplier} tables, which are joined
repeatedly by the 22 queries and whose aggregate size is far below 4\,GiB, so
they remain resident; the one-shot large scan of \texttt{lineitem} is correctly
refused admission by the frequency filter and causes no scan pollution.

\subsubsection{Effectiveness of Layout Reconfiguration: Fragmentation of the Canonical Layout}

The effectiveness of layout reconfiguration is governed by the amount of
coalesceable fragmentation present in the canonical layout. Measuring that
fragmentation as the number of remote requests per GiB of data yields
Table~\ref{tab:fragmentation}.

\begin{table}[t]
\centering
\caption{Fragmentation of the canonical layout and the gain achieved by layout
reconfiguration alone.}
\label{tab:fragmentation}
\footnotesize
\setlength{\tabcolsep}{4pt}
\begin{tabular}{@{}lrrr@{}}
\toprule
Dataset & \makecell[r]{Fragmentation\\(GET / GiB)} &
\makecell[r]{GET reduction\\by layout} & \makecell[r]{Wall-clock gain\\(AWS)} \\
\midrule
TPC-H SF350      & 11{,}422 & 76.4\% & 26.3\% \\
NYC Yellow Taxi  & 3{,}099  & 92.4\% & 17.6\% \\
ClickBench SF8   & 2{,}342  & 30.0\% & 22.8\% \\
Criteo CTR       & 1{,}869  & 45.7\% & 19.0\% \\
NovoMCP Compound & 634      & 22.5\% & \textbf{3.8\%} \\
\bottomrule
\end{tabular}
\end{table}

This relationship is likewise of threshold form rather than monotonic: the four
datasets whose fragmentation is at least 1{,}869\,GET/GiB all attain wall-clock
gains between 17.6\% and 26.3\%, whereas NovoMCP, with the lowest
fragmentation, attains only 3.8\%. The reason for NovoMCP's ineffectiveness is
structural: its canonical layout already comprises only 21 files and 132 row
groups, so that \textbf{it is already close to the optimal geometry for this
workload, and reconfiguration can reduce the file count only from 21 to 20,
leaving no margin for coalescing}. In other words, the gain from layout
reconfiguration derives from the gap between the canonical and the optimal
layout, and where that gap is itself small, the mechanism has no room to act.

Above the threshold, the gain does not increase monotonically with
fragmentation, because the dimension along which coalescing is possible differs
across datasets: the request reduction on NYC arises chiefly from row-group
coalescing (its 204 canonical files contain 14{,}973 row groups), whereas on
Criteo it arises chiefly from file coalescing (973 files merged into 32).
Fragmentation thus bounds the space available to the mechanism, while the
fraction actually realized depends on whether the dimension in which the
fragmentation resides lies within the action space.

\textbf{The two conditions are mutually independent.}
The capacity pressure ratio characterizes a temporal property of the queries,
whereas fragmentation characterizes a spatial property of the data, and the two
do not covary within our sample: TPC-H has both the highest pressure ratio and
the highest fragmentation, whereas NovoMCP conversely has a pressure ratio of
0.7$\times$ together with the lowest fragmentation. No single mechanism can
therefore cover the full range of workloads, while the combination of the two
has at least one effective member on each of the five datasets studied here.
This constitutes the direct justification for adopting a combined design rather
than a single mechanism.

\subsubsection{Environment Sensitivity: Latency Amplification and the Share of Request Reduction}

The gain of a given mechanism is not amplified proportionally across the two
cloud environments; its amplification factor is governed by the \emph{share of
request reduction} in that mechanism's gain. Comparing the ratio of the
wall-clock gain of caching on COS to that on AWS against its GET reduction
yields a Spearman rank correlation of 0.80.

\begin{table}[t]
\centering
\caption{Amplification of the caching gain from AWS to COS.}
\label{tab:amplification}
\footnotesize
\begin{tabular}{@{}lrr@{}}
\toprule
Dataset & \makecell[r]{GET reduction\\by caching} &
\makecell[r]{Amplification of\\the caching gain} \\
\midrule
NYC Yellow Taxi  & 10.6\% & 1.58 \\
TPC-H SF350      & 11.3\% & 1.10 \\
ClickBench SF8   & 73.0\% & 1.99 \\
Criteo CTR       & 74.8\% & 1.59 \\
NovoMCP Compound & 77.2\% & 3.10 \\
\bottomrule
\end{tabular}
\end{table}

The mechanism is consistent with the conclusion of \S\ref{subsec:main}: the
fixed per-request cost rises from 26.1\,ms on AWS to 86.2\,ms on COS while
bandwidth is not the bottleneck, so that \textbf{the portion of a mechanism's
gain contributed by request reduction is amplified in proportion to the
environment latency, whereas the portion contributed by byte reduction is not}.
Caching reduces only requests, while layout reconfiguration reduces both
requests and bytes; the amplification of caching in a high-latency environment
is therefore generally larger than that of layout reconfiguration.

This regularity accounts for the reversal in the ordering of the two mechanisms
on ClickBench SF8. This dataset is the one in our sample where the request
reduction of caching exceeds that of layout reconfiguration by the largest
margin, so that on migrating from AWS to COS the caching gain is amplified by a
factor of 1.99 while the layout gain is scarcely amplified at all; the ordering
is thereby reversed from ``layout 22.8\% versus caching 18.3\%'' on AWS to
``caching 36.5\% versus layout 23.0\%'' on COS. Criteo exhibits a shift in the
same direction that stops short of a reversal: on AWS layout leads by 19.0\% to
13.1\%, while on COS the two converge to 21.4\% and 20.8\%.

Our quantitative evidence for environment amplification rests on Criteo, whose
I/O metrics agree almost digit for digit across the two clouds (see
\S\ref{subsec:protocol}).

\subsubsection{Scale Sensitivity}
\label{subsec:scale}

With the cache capacity fixed at 4\,GiB, growth in data scale acts in two
opposing directions. Taking ClickBench from SF1 (14.11\,GiB) to SF8
(111.99\,GiB) on AWS:

\begin{table}[t]
\centering
\caption{Effect of data scale on the two mechanisms (ClickBench, AWS).}
\label{tab:scale}
\footnotesize
\begin{tabular}{@{}lrr@{}}
\toprule
Effect & SF1 & SF8 \\
\midrule
GET reduction by caching   & 85.5\% & 73.0\% \\
Wall-clock gain of caching & 14.9\% & 18.3\% \\
GET reduction by layout    & 29.1\% & 30.0\% \\
Wall-clock gain of layout  & 14.5\% & 22.8\% \\
Combined wall-clock gain   & 21.7\% & 31.6\% \\
\bottomrule
\end{tabular}
\end{table}

Growth in scale raises the total remote read volume of the canonical baseline
from 46.6 to 371.9\,GiB and the mean per-query read volume from 1.08 to
8.65\,GiB, so the request reduction of caching falls from 85.5\% to 73.0\%,
consistent with the relationship of \S\ref{subsec:pressure}. The wall-clock
gains of both mechanisms nevertheless rise with scale, and the combined gain
increases from 21.7\% to 31.6\%. The reason is that \textbf{the share of remote
reads in wall-clock time grows with scale}: the same proportional physical
reduction is realized as a larger wall-clock gain on a workload with a higher
I/O share, and this effect outweighs the decline in the proportional reduction
itself. No such offset arises on the layout side, whose proportional request
reduction is nearly invariant with scale, so its wall-clock gain grows more
markedly. One corollary concerning capacity provisioning follows: with the cache
capacity held fixed, growth in data scale increases the share of layout
reconfiguration in the combined gain.

\subsubsection{Interaction Between the Two Mechanisms}

In eight of the ten grids the combined gain is smaller than the sum of the two
individual gains, i.e., the interaction term is negative. The mechanism is that
the two compete for the same portion of the gain: in coalescing physical reads,
layout reconfiguration has already eliminated some of the small requests that
the cache would otherwise have served as hits. This overlap can be quantified
directly---on ClickBench SF8 / AWS, caching eliminates 191{,}465 requests on the
canonical layout but only 139{,}078 on the reconfigured layout, the difference of
52{,}387 being the portion already absorbed by the layout.

The two TPC-H SF350 grids are the exception, with a positive interaction term.
The mechanism is the converse of the overlap above: reconfiguration contracts
the file count from 2{,}130 to 514 and thereby shortens reuse distances
substantially, bringing a working set that previously could not fit within the
4\,GiB capacity into the effective regime of the cache; the hit rate rises from
12.3\% to 59.1\% and the absolute number of cache hits increases. This is
consistent with the conclusion of \S\ref{subsec:pressure}: layout
reconfiguration lowers the capacity pressure ratio and thereby moves the cache
from its ineffective into its effective regime.

The origin of this positive interaction can be read directly from the paired
data: the per-round gain of caching on the reconfigured layout is stable between
261.1 and 317.2\,s, whereas on the canonical layout it decays from 268.3\,s in
the first round to 91.3\,s in the fifth. Under the canonical layout the capacity
pressure ratio is 6.7 and whether a hit occurs depends on incidental runtime
factors; once reconfiguration brings the ratio into the effective regime, the
caching gain becomes deterministic. \textbf{Layout reconfiguration therefore not
only raises the mean of the caching gain but also reduces its variance.}

\subsection{Cache Algorithm Selection}
\label{subsec:policy}

This section addresses why the cache tier adopts cost-aware W-TinyLFU. We
compare it against five cache algorithms and a no-cache baseline on a single
workload: the reconfigured layout of ClickBench SF8, with the capacity fixed at
4\,GiB, Spark \texttt{local[16]}, and a driver heap of 32\,GiB, with $n=5$ for
each of the seven arms. The arms differ only in the eviction and admission
policy; all other configuration is identical.

\subsubsection{Choice of Baseline Algorithms}

From five mainstream families of cache policy---pure recency, segmented
recency, FIFO-based eviction, frequency-based admission, and size-aware
admission---we select one representative or recently proposed algorithm each as
a baseline, and evaluate them together with the no-cache baseline under the same
4\,GiB memory budget. Specifically: (1)~\emph{LRU}~\cite{lru}, the classical
baseline for pure-recency eviction, which exhibits cache behaviour in the
absence of scan protection; (2)~\emph{SLRU}~\cite{slru} (with the protected
segment occupying 80\% of the capacity), representing segmented-recency
eviction, which isolates one-shot large scans by means of a probationary and a
protected segment and corresponds to the structure of our main cache region once
the admission filter is removed; (3)~\emph{SIEVE}~\cite{sieve}, representing the
recently proposed family of minimal FIFO-based eviction algorithms, which
realizes lazy promotion and quick demotion through a single scanning pointer and
serves as a state-of-the-art low-overhead eviction baseline;
(4)~\emph{W-TinyLFU (plain)}~\cite{wtinylfu}, representing admission control
based on Count--Min-Sketch frequency estimation, which makes admission decisions
on the frequency of individual blocks; and (5)~\emph{W-TinyLFU-AV}~\cite{einziger2022av},
representing the size-aware extension of W-TinyLFU for variable-length objects,
which compares the frequency of a new object against the aggregate frequency of
the set of prospective victims and is the closest recent work to our cost-aware
admission.

\subsubsection{Results of the Seven-Arm Matrix}

\begin{table}[t]
\centering
\caption{Cache algorithm matrix on the reconfigured layout of ClickBench SF8,
4\,GiB capacity, $n=5$ per arm. The coefficient of variation lies between
0.23\% and 1.38\% across the seven arms, well below our 5\% threshold. Remote
bytes, hit counts, GC time, and evicted volume are omitted here and discussed
in the text.}
\label{tab:policy}
\footnotesize
\setlength{\tabcolsep}{5pt}
\begin{tabular}{@{}lrrrr@{}}
\toprule
Algorithm & \makecell[r]{Median\\(s)} & \makecell[r]{vs.\\no cache} &
GET & \makecell[r]{Peak heap\\(GiB)} \\
\midrule
SIEVE                         & \textbf{1651.7} & 12.37\% & 45{,}943 & 27.51 \\
\textbf{Cost-aware W-TinyLFU} & 1656.1 & 12.14\% & \textbf{44{,}116} & \textbf{23.90} \\
W-TinyLFU-AV                  & 1664.4 & 11.69\% & 48{,}481 & 28.61 \\
W-TinyLFU (plain)             & 1677.1 & 11.02\% & 49{,}519 & 28.63 \\
SLRU                          & 1700.8 & 9.76\%  & 64{,}259 & 28.31 \\
LRU                           & 1731.9 & 8.11\%  & 86{,}250 & 28.03 \\
No cache                      & 1884.8 & ---     & 183{,}675 & 12.42 \\
\bottomrule
\end{tabular}
\end{table}

The wall-clock ordering of the seven arms coincides exactly with their ordering
by request reduction, which again corroborates the conclusion of
\S\ref{subsec:main}. The spread between the extremes quantifies the two axes of
our design: cost-aware W-TinyLFU issues 10.9\% fewer requests than plain
W-TinyLFU, 31.4\% fewer than SLRU---which is our main region with the admission
filter removed---and 48.8\% fewer than pure-recency LRU.

\subsubsection{Basis for Selection: Equal Wall Clock at Substantially Lower Cost}

Our scheme and SIEVE, the fastest by median, are \textbf{indistinguishable} in
wall clock, differing by only 0.26\%, and the win-loss ratio for each round is 3 to 2. According to the criteria in Section 5.1.3, there is no winner. The selection therefore rests on the cost
incurred to reach that same wall clock, and our scheme is clearly superior on
four cost metrics: \textbf{its peak heap is 3.61\,GiB (13\%) lower, it issues
1{,}827 fewer requests, it spends 1.68\,s less in garbage collection, and its
evicted volume is two orders of magnitude smaller (0.44 versus 64.13\,GiB)}. Its
hit count is at the same time the highest in the matrix (139{,}559, against
137{,}732 for SIEVE).

Two mechanisms account for this difference in cost.

\textbf{Request-level atomic admission lets the same volume of bytes cover fewer
requests.}
Our scheme and SIEVE read almost identical remote volumes (210.65 versus
210.64\,GiB), yet ours issues 1{,}827 fewer requests, so that the mean size of a
request rises from 4.69 to 4.89\,MiB. The reason is that we take a single Range
GET as the atomic unit of admission: what the cache retains is a complete
logical unit of reuse, and one hit therefore eliminates one entire remote
request. Per-block admission, by contrast, leaves partially covered intervals,
so a subsequent read must still issue a request to supply the missing blocks;
bytes are saved while the request is not eliminated. This result is consistent
with the design rationale for request-level atomic admission given in
\S\ref{subsec:cachealg}.

\textbf{The cost-aware admission criterion drives cache turnover towards zero.}
Our scheme admits a candidate request only if its benefit strictly exceeds the
sum of the benefits of all victims it would displace. Once the cache has been
filled with high-benefit content, subsequent candidates rarely clear that
threshold, so both admissions and evictions fall sharply: the evicted volume is
a mere 0.44\,GiB, whereas all six other algorithms exceed 64\,GiB. Low turnover
directly reduces the volume of transient allocation, which accounts jointly for
the lower peak heap and the shortest garbage-collection time.

Because the scalability ceiling of an in-memory cache tier is set by the heap
rather than by the hit rate (see \S\ref{subsec:capacity}), a lower heap
footprint at equal wall clock means more usable capacity, or higher
per-host concurrency, under the same resources. We therefore adopt cost-aware
W-TinyLFU.

\subsection{Cache Capacity Curves and the Limits of the In-Memory Tier}
\label{subsec:capacity}

The main experiments fix the cache capacity at 4\,GiB in order to keep the
comparison across workloads fair. This section gives the basis for that choice
and answers the two questions that follow from it: where the returns from
additional capacity are exhausted, and what should be done when memory cannot
supply the capacity required.

\subsubsection{Capacity--Benefit Curves at Different Pressure Ratios}

We measure capacity curves on two workloads whose pressure ratios differ by an
order of magnitude. NYC Yellow Taxi has a pressure ratio of 6.5 at 4\,GiB and
thus lies on the ineffective side of the cliff described in
\S\ref{subsec:pressure}, whereas ClickBench SF8, at 0.68, lies on the effective
side. We choose NYC Yellow rather than TPC-H because layout reconfiguration does
not reduce the volume of bytes accessed on NYC Yellow and therefore cannot lower
its capacity pressure ratio, unlike the situation on TPC-H.

The sweep over six arms on the canonical layout of NYC in the COS environment is
shown in Fig.~\ref{fig:cap_nyc}. The step gain from each doubling of capacity
\emph{increases}. Relative to the 4\,GiB reference arm, the same-round paired
differences on the 8, 16, and 32\,GiB arms correspond to relative margins of
15.0\%, 33.4\%, and 57.4\%, respectively. The complete working set of a single
epoch is measured at 25.98\,GiB and therefore does not fill the 32\,GiB cache,
so that arm's effective pressure ratio is exactly 1.0. On that arm the
\textbf{number of rejected blocks, of evicted blocks, and of ghost hits are all
zero}, and the useful bytes supplied by the cache reach the asymptotic upper
bound for four epochs (one cold scan followed by three hot hits); the request
count falls to 30{,}068 ($-75.3\%$) and remote bytes to 25.99\,GiB ($-75.0\%$),
both equal to the lower bound at which every byte is fetched from the remote
store exactly once. Further enlarging the capacity therefore cannot produce a
single additional hit: the plateau has been reached.

The four-arm sweep on ClickBench SF8 is shown in Fig.~\ref{fig:cap_cb}. The gain
is concentrated at the 4\,GiB cache ($-14.5\%$); thereafter 4$\to$8\,GiB yields
only a further 2.6\% and 8$\to$16\,GiB yields $-0.38\%$. The reason for the
diminishing cache gain is evident in the eviction data: on both the 8 and the
16\,GiB arms the evicted volume over a full round is 0.00\,GiB, and the number of
blocks requested again after eviction is 1 and 0, respectively. The sweep on
ClickBench SF1 further exhibits working-set capping---the resident cache on both
the 4 and the 16\,GiB arms stops at 2.67\,GiB, with identical request counts
(4{,}798) and remote bytes (36.2\,GiB) digit for digit.

Arranging the ten arms by pressure ratio shows that \textbf{the knee lies at a
pressure ratio of exactly 1}, that is, at the point where capacity just covers
the working set of one pass. This relationship establishes the capacity setting
used in the main experiments: three of our five workloads (ClickBench, Criteo,
NovoMCP) have pressure ratios below 1.2, for which 4\,GiB already places them at
the plateau of the benefit curve; of the remaining two, TPC-H can have its
pressure ratio lowered by layout reconfiguration while NYC cannot. 4\,GiB is
therefore a reasonable setting that covers the majority of workloads rather than
the optimal configuration for all of them---\textbf{our main results consequently
understate the potential of caching on workloads with higher pressure ratios}.

\subsubsection{Local SSD as Off-Heap Capacity: Feasibility}

Given that additional capacity yields substantial gains on workloads with high
pressure ratios, the question that follows is whether that capacity can be
supplied by the JVM heap. Our measured marginal heap cost is 2.85 to 3.80 times
for the step from no cache to 4\,GiB, converging to 1.09 to 1.19 times for every
step above 8\,GiB. The higher multiplier of the first step stems chiefly from the
one-time overhead of enabling the cache---SDK and Spark objects kept alive along
the hit path, tee buffers, and the like---which does not scale linearly with
capacity. Experiments show that with a 32 GiB cache, the peak heap size is 40.49 GiB, which has triggered a heap water level warning. Once the working set grows to the scale
of hundreds of GiB, or once a single host must retain a working set for each of
several concurrent tenants, the capacity required exceeds any single-host heap
configuration.

The capacity limit of the in-memory cache implies that the portion of capacity
beyond the heap watermark must be carried by another medium. A local SSD is the
most immediate candidate. Its viability presupposes that its hit latency be
markedly lower than the GET latency of the remote object store. We measure two
local media on an AWS \texttt{m5dn.4xlarge} using \texttt{fio}. At the 64 to
256\,KiB block sizes that represent the granularity of a cache hit, the p99
latency of instance-store NVMe is 0.18 to 0.51\,ms, approximately \textbf{1/46}
of the p50 of a small same-region S3 request, so the precondition holds. On
workloads with high pressure ratios it is therefore feasible, on the medium side,
to push the capacity beyond the heap watermark down to local NVMe. We accordingly
leave tiered caching to future work.

\subsubsection{Two Lines of Future Work Implied by the Capacity Curves}

\textbf{Adaptive capacity provisioning.}
The main experiments fix the capacity at 4\,GiB for the sake of fair comparison.
That setting places workloads with a pressure ratio below 1 at the plateau of
the benefit curve, but realizes only a small fraction of the potential on
workloads such as NYC (a 16.0\% wall-clock reduction at 4\,GiB, against 63.9\% at
26\,GiB). Since the working set of a single pass can be estimated online from the
request-level traces we already collect, and since the benefit is approximately
linear in the ratio of capacity to working set, capacity admits a basis for
automatic per-workload provisioning: the objective is to drive the pressure ratio
towards 1 while halting growth as the heap watermark approaches its limit. Our
ghost table already supplies the required signal---the number of blocks
``requested again after rejection'' is non-zero at 4\,GiB and falls to zero at
32\,GiB, which can serve as direct feedback on whether capacity is sufficient.

\textbf{Off-heap and local SSD tiering.}
For larger working sets, memory cannot supply a sufficiently large cache.
\S\ref{subsec:ssd} establishes the preliminary feasibility of such tiering.
\subsection{Action Ablation and Planner Fidelity}
\label{subsec:ablation}

This subsection identifies the sources of the gain delivered by layout
reconfiguration. We first quantify the end-to-end contribution of each action
through a leave-one-out ablation, then add positive controls for the actions
whose signal is weak in order to distinguish an ineffective action from one that
is insufficiently exercised by the workload, and finally validate the ranking
fidelity of the planner.

\subsubsection{Leave-One-Out Ablation: Row Groups and Compression Dominate}

On TPC-H SF100 we remove one class of action at a time from the full plan while
keeping all others unchanged; the results are given in
Table~\ref{tab:ablation}.

\begin{table}[t]
\centering
\caption{Leave-one-out action ablation on TPC-H SF100.}
\label{tab:ablation}
\footnotesize
\setlength{\tabcolsep}{4pt}
\begin{tabular}{@{}lrrrr@{}}
\toprule
Arm & \makecell[r]{Median\\(s)} & \makecell[r]{vs.\\full plan} & GET &
\makecell[r]{Remote\\bytes (GB)} \\
\midrule
Canonical layout       & 2{,}688.494 & +112.0\% & 383{,}787 & 171.7 \\
Full plan              & 1{,}267.958 & ---      & 97{,}967  & 117.7 \\
\midrule
$-$\,row-group         & 4{,}427.962 & \textbf{+249.2\%} & 694{,}107 & 191.6 \\
$-$\,compression       & 1{,}320.058 & \textbf{+4.11\%}  & 100{,}661 & 166.5 \\
$-$\,encoding          & 1{,}274.748 & +0.54\%  & 98{,}093  & 118.9 \\
$-$\,file-size         & 1{,}273.683 & +0.45\%  & 97{,}033  & 117.5 \\
$-$\,column-order      & 1{,}272.760 & +0.38\%  & 100{,}568 & 117.7 \\
\bottomrule
\end{tabular}
\end{table}

The data show that on the TPC-H workload, the row-group action has by far the
largest effect on layout optimization; the layout with row-group restructuring
removed is even markedly worse than the canonical layout. The reason is that the
canonical \texttt{lineitem}, \texttt{orders}, and \texttt{partsupp} tables
contain a large number of small row groups of only 8--17\,MiB; each row group
yields its own chunk for every projected column and frequently triggers the
64\,KiB minimum range read. Enlarging the row groups of the main tables to
roughly 70--117\,MiB substantially reduces both the number of column chunks and
the number of range GETs. \textbf{Compression ranks second}, and is the only
action in the sample whose gain arises chiefly from reducing bytes rather than
requests.

The signal of the remaining three actions on this workload is below 1\%. This
does not indicate that they are ineffective in general; it is a consequence of
\textbf{unbalanced action exposure}. The full plan for TPC-H SF100 contains only
four explicit encoding actions, all of which fall on the \texttt{partsupp}
table, whereas the two main tables \texttt{lineitem} and \texttt{orders}, which
account for more than 70\% of the I/O of the whole database, carry no encoding
action at all; column order likewise applies only to tables outside the main
scan. Under such a distribution the leave-one-out ablation yields only a lower
bound, and we therefore add positive controls for these three actions
(\S\ref{subsec:posctrl}).

\subsubsection{Positive Controls for Weak-Signal Actions}
\label{subsec:posctrl}

A positive control constructs two layouts that differ only in the target action
and in which that action is fully exercised, in order to test whether the action
itself is capable of producing a gain.

\textbf{File count (Criteo, $n=5$).}

\begin{table}[t]
\centering
\caption{Positive control for target file size (Criteo).}
\label{tab:posctrl_file}
\footnotesize
\begin{tabular}{@{}lrrr@{}}
\toprule
Arm & Median (s) & GET & vs.\ full plan \\
\midrule
Full plan          & 1{,}074.858 & 44{,}242 & ---      \\
$-$\,file-size     & 1{,}217.473 & 71{,}762 & +13.3\%  \\
\bottomrule
\end{tabular}
\end{table}

Removing the action adds 142.6\,s of wall clock and increases the request count
by 62.2\%, while remote bytes actually fall by 0.46\%. The gain of this action
therefore arises entirely from the object count and the request overhead that
accompanies it, consistent with the conclusion of \S\ref{subsec:main}.

\textbf{Encoding and column order (ClickBench, $n=5$).}

\begin{table}[t]
\centering
\caption{Positive controls for encoding and column order (ClickBench).}
\label{tab:posctrl_cb}
\footnotesize
\begin{tabular}{@{}lrrr@{}}
\toprule
Arm & Median (s) & vs.\ full plan & GET \\
\midrule
Full plan           & 321.900 & ---        & 23{,}588 \\
$-$\,encoding       & 336.837 & +4.64\%    & 24{,}192 \\
$-$\,column-order   & 321.104 & $-$0.25\%  & 25{,}116 \\
\bottomrule
\end{tabular}
\end{table}

With encoding removed the workload is slower in all five rounds, which
constitutes a conservative positive case showing that encoding does yield an
end-to-end gain when adequately exercised. Column order, by contrast, reduces the
request count by 6.5\% while the end-to-end direction is mixed. The physical
mechanism behind the absence of a clear effect for this action is analysed
separately in \S\ref{subsec:colorder}.

\textbf{Encoding (Criteo, $n=5$).}
To test whether the encoding conclusion above depends on the particular
workload, we repeat this positive control on Criteo.

\begin{table}[t]
\centering
\caption{Positive control for encoding (Criteo).}
\label{tab:posctrl_enc}
\footnotesize
\begin{tabular}{@{}lrrr@{}}
\toprule
Arm & Median (s) & vs.\ full plan & GET \\
\midrule
Full plan       & 1{,}093.60 & ---      & 38{,}048 \\
$-$\,encoding   & 1{,}143.45 & +4.56\%  & 38{,}888 \\
\bottomrule
\end{tabular}
\end{table}

On ClickBench, removing encoding actually produced 21 files rather than the 22
of the full plan. On Criteo, by contrast, the file and row-group geometry of the
two arms is identical item by item, so the difference can only arise from the
encoding of those four columns and not from a change in file splitting. The
three encoding results jointly support the exposure explanation of the
preceding subsection: when encoding acts on the main scan path it yields a small
but stable positive gain, whereas when it acts only on tables outside the main
scan (\texttt{partsupp} in TPC-H), the whole database shows no end-to-end gain.

\subsubsection{Column Order: Request Reduction and Concurrency on the Critical Path}
\label{subsec:colorder}

In the leave-one-out ablation, column order consistently reduces the number of
remote GET requests, yet under the default read configuration it yields no
significant reduction in end-to-end latency. To locate the cause, we construct
on ClickBench two control geometry levels with strictly fixed file and
row-group boundaries: (1)~\texttt{production-RG} (22 files, 169 row groups),
corresponding to the physical layout at a conventional row-group size; and
(2)~\texttt{small-RG} (73 files, 578 row groups, with a 64\,MiB row-group
target), which shrinks the row groups and the individual column chunks so that
more adjacent columns fall within the 2\,MiB coalescing window of S3A vectored
reads, thereby amplifying the effect of column order on request coalescing. At
each level we compare a \texttt{co-located} layout, in which co-accessed hot
columns are placed contiguously, against a \texttt{separated} layout, in which
cold columns are interleaved between the hot ones.

Our hypothesis is that column order fails to shorten end-to-end wall-clock time
because the requests it eliminates lie inside the concurrent read batch of a
single row group rather than on the critical path of the query. Specifically,
when the Parquet reader opens a row group it submits the byte ranges of all
projected columns in that row group to the \texttt{readVectored} interface at
once, and S3A issues the range requests in parallel up to a per-stream
concurrency limit (4 by default). Let $\mathrm{Union}$ denote the length of the
union of the request time intervals (i.e., the critical-path window actually
occupied by I/O) and $K_{\mathrm{busy}}$ the mean number of in-flight requests
during busy periods; the sum of the time-to-first-byte of all requests,
$\sum \mathrm{TTFB}$, then satisfies the identity
$$
  \sum \mathrm{TTFB} = \mathrm{Union} \times K_{\mathrm{busy}}.
$$

Under the \texttt{separated} layout, columns that would otherwise be adjacent are
split into multiple independent range requests, but these requests are in
flight concurrently within the same read stream, so the I/O wait of the row
group is determined by the slowest request in the batch rather than by the sum
of all request latencies. Rearranging the hot columns into the
\texttt{co-located} layout coalesces physically adjacent column chunks into
fewer GET requests and directly lowers $\sum \mathrm{TTFB}$, but it lowers the
in-flight concurrency $K_{\mathrm{busy}}$ correspondingly, so that the
critical-path I/O window $\mathrm{Union}$ remains almost unchanged.

Request-level traces corroborate this relationship directly. For the lightweight
projection queries on \texttt{small-RG}, the \texttt{co-located} layout reduces
the number of GET requests by 44.3\% relative to \texttt{separated}, and
$\sum \mathrm{TTFB}$ falls from 619.2\,s to 385.7\,s ($-37.7\%$); at the same
time, however, $K_{\mathrm{busy}}$ falls in step from 21.41 to 13.52
($-36.9\%$), so that the critical-path window $\mathrm{Union}$ shrinks only from
30.17\,s to 27.95\,s ($-7.4\%$). Across the full set of 43 queries, column order
lowers $K_{\mathrm{busy}}$ from 11.92 to 11.06, while $\mathrm{Union}$ differs by
only 0.07\% (511.06\,s versus 510.68\,s), and the end-to-end difference falls
entirely within run-to-run noise.

To further verify that per-stream concurrency is the direct cause masking the
gain of column order, we lower the S3A per-stream vectored concurrency limit from
its default of 4 to 1 and repeat the measurements on both geometry levels under
two degrees of task parallelism, 16 threads and a single thread.

\begin{table}[t]
\centering
\caption{Effect of the per-stream vectored concurrency limit on the gain from
column order (ClickBench).}
\label{tab:colorder_knob}
\footnotesize
\setlength{\tabcolsep}{3.5pt}
\begin{tabular}{@{}llrrr@{}}
\toprule
Configuration & Level & \makecell[r]{\texttt{co-loc.}\\(s)} &
\makecell[r]{\texttt{sep.}\\(s)} & Gain \\
\midrule
\texttt{local[16]}, \texttt{active=4}
  & \texttt{pd-RG} & 98.79  & 99.96  & $-1.2\%$ \\
  & \texttt{small-RG}      & 170.89 & 174.97 & $-2.3\%$ \\
\addlinespace[2pt]
\texttt{local[16]}, \textbf{\texttt{active=1}}
  & \texttt{pd-RG} & 109.66 & 115.51 & $\mathbf{-5.1\%}$ (3/3) \\
  & \texttt{small-RG}      & 197.07 & 234.23 & $\mathbf{-15.9\%}$ (3/3) \\
\addlinespace[2pt]
\texttt{local[1]}, \textbf{\texttt{active=1}}
  & \texttt{pdn-RG} & 313.91 & 358.46 & $\mathbf{-12.4\%}$ (3/3) \\
  & \texttt{small-RG}      & 585.66 & 875.72 & $\mathbf{-33.1\%}$ (3/3) \\
\bottomrule
\end{tabular}
\end{table}

The results show that as soon as the number of concurrent requests within a
single stream is limited to 1, so that the range requests split within a row
group degrade into serial waits, the end-to-end gain of column order appears
immediately and is consistent in direction across all three rounds: under
\texttt{local[16]}, \texttt{production-RG} and \texttt{small-RG} achieve
wall-clock speedups of 5.1\% and 15.9\%, respectively; under the fully serial
\texttt{local[1]}, the wall-clock time of \texttt{small-RG} falls from 875.72\,s
to 585.66\,s ($-33.1\%$), and the 290.1\,s of end-to-end time saved matches the
280\,s difference in cumulative request latency almost one to one. These results
explain mechanistically the conclusion of \S\ref{subsec:main} that end-to-end
time is determined by request waits on the critical path: whether a reduction in
remote GET requests translates into a wall-clock gain depends on whether the
eliminated requests are serial waits or concurrent overlaps on the execution
path. Adjusting file size and row-group size eliminates cross-file and
cross-row-group scheduling and metadata requests, and online caching eliminates
repeated requests across queries and training epochs; both classes lie on the
serial critical path. Column order, by contrast, eliminates range requests
inside the \texttt{readVectored} batch of a single row group, which under the
default concurrency configuration are already masked by multi-request
concurrency within a stream. Column order is therefore a secondary optimization
dimension constrained by the concurrency of the underlying reader; we retain this
dimension in the candidate space to accommodate different client concurrency
configurations, but do not treat it as a principal source of end-to-end speedup
on the default high-concurrency read path.

\subsubsection{Ranking Fidelity and Geometric Closure of the Planner}
\label{subsec:fidelity}

The action ablation establishes the contribution of each action; the remaining
question is whether the planner chooses correctly among candidates. On Criteo we
freeze 8 of the 272 candidates that satisfy the feasibility constraints, rank
them by L1 solely on their predicted I/O time, and then materialize and measure
all of them with $n=5$.

\begin{table}[t]
\centering
\caption{Ranking fidelity and geometric prediction error of L1 (Criteo, 8
candidates).}
\label{tab:fidelity}
\footnotesize
\begin{tabular}{@{}lr@{}}
\toprule
Metric & Value \\
\midrule
Spearman rank correlation                         & \textbf{0.9519} \\
Kendall $\tau$-b                                  & 0.8895 \\
Selection regret                                  & 2.15\% \\
Row-group count error (median / max)              & 0.00\% / 0.00\% \\
File count error (median / max)                   & 0.00\% / 10.25\% \\
Compressed-byte error (median / max)              & 0.67\% / 2.95\% \\
\bottomrule
\end{tabular}
\end{table}

L1 is therefore an accurate relative ranker, and its underlying virtual footer
achieves tight geometric closure with the physically written output across the 8
candidates: the row-group counts match exactly in every case, and the maximum
error in compressed size is below 3\%. This closure holds equally for the
production reconfiguration plans of all five datasets (six scales), each of which
passes, before upload, the constraint of a row-group count error of at most one
and file-count and compressed-byte errors each of at most 15\%.
